% Part: propositional-logic
% Chapter: syntax-and-semantics
% Section: introduction

\documentclass[../../../include/open-logic-section]{subfiles}

\begin{document}

\olfileid{pl}{syn}{int}

\olsection{ভূমিকা}

বচনমূলক যুক্তিবিদ্যা সেই সব !!{formula}s নিয়ে আলোচনা করে, যেগুলি
!!{propositional variable}s থেকে $\lnot$, $\land$, $\lor$, $\lif$ ও
$\liff$ বচনসংযোজকের সাহায্যে তৈরি হয়। স্বজ্ঞাগতভাবে,
!!a{propositional variable}~$p$ এমন কোনো বাক্য বা বচনের স্থলাভিষিক্ত,
যা সত্য অথবা মিথ্যা। !!a{formula}-এর !!{propositional variable}-গুলির
“সত্যমান” নির্ধারিত হলেই, সেগুলি থেকে বচনসংযোজকের সাহায্যে তৈরি যেকোনো
!!{formula}s-এর সত্যমানও নির্ধারিত হয়। বচনমূলক যুক্তিবিদ্যাকে আমরা
\emph{সত্যমান-অপেক্ষকধর্মী} বলি, কারণ এর অর্থতত্ত্ব সত্যমানের অপেক্ষকের
সাহায্যে দেওয়া হয়। বিশেষত, বচনমূলক যুক্তিবিদ্যায় সত্য বা মিথ্যা হওয়ার
আরও সূক্ষ্ম নির্ধারণগুলি আমরা বিবেচনা করি না; যেমন, কোনো বিষয় নিছক আপতিকভাবে
সত্য না হয়ে আবশ্যিকভাবে সত্য কি না, তা সত্য বলে জানা আছে কি না, অথবা তা এখন
সত্য—আগে সত্য ছিল কিংবা পরে সত্য হবে কি না। আমরা শুধু সত্য ($\True$) ও
মিথ্যা ($\False$)—এই দুই সত্যমান বিবেচনা করি। তাই কোনো উক্তি সত্যও নয়,
মিথ্যাও নয়, অথবা অর্ধসত্য—এমন সম্ভাবনা আলোচনা থেকে বাদ দিই। আমরা শুধু সেই
সব সংযোজকের উপরও মন দিই, যেগুলি দিয়ে তৈরি !!a{formula}-এর সত্যমান তার
অংশগুলির সত্যমান দিয়ে সম্পূর্ণ নির্ধারিত হয় (ধরা যাক, তার অর্থ দিয়ে নয়)।
বিশেষত, ইংরেজি ভাষার শর্তবচনের সত্যমান এই অর্থে সত্যমান-অপেক্ষকধর্মী কি না,
তা বিতর্কিত। বস্তুগত শর্তবচন~$\lif$ সত্যমান-অপেক্ষকধর্মী; অন্য যুক্তিবিদ্যায়
এমন শর্তবচন নিয়ে আলোচনা করা হয়, যেগুলি সত্যমান-অপেক্ষকধর্মী নয়।

সত্যমান-অপেক্ষক বচনমূলক যুক্তিবিদ্যার তত্ত্ব ও অধিতত্ত্ব গড়ে তুলতে প্রথমে
তার অভিব্যক্তিগুলির সংকেতবিন্যাস ও অর্থতত্ত্ব সংজ্ঞায়িত করতে হবে।
সংযোজক ব্যবহার করে !!{propositional variable}s থেকে !!{formula}s তৈরির
একটি উপায় আমরা বর্ণনা করব। বিকল্প সংজ্ঞাও সম্ভব। অন্য পদ্ধতিগুলি আলাদা
সংকেত বেছে নিতে পারে, আলাদা সংযোজকসমষ্টিকে মৌলিক ধরতে পারে এবং বন্ধনী
আলাদাভাবে ব্যবহার করতে পারে (অথবা তথাকথিত পোলিশ সংকেতরীতির মতো একেবারেই
ব্যবহার না-ও করতে পারে)। তবে সব পদ্ধতির একটি সাধারণ বৈশিষ্ট্য হলো:
গঠন-বিধিগুলি !!{formula}s-এর সেটকে \emph{আরোহীভাবে} সংজ্ঞায়িত করে।
কাজটি ঠিকভাবে করা হলে গঠন-বিধি অনুসারে প্রতিটি অভিব্যক্তি মূলত একটিমাত্র
উপায়েই তৈরি হতে পারে। যে আরোহী সংজ্ঞা \emph{এককভাবে পাঠযোগ্য}
অভিব্যক্তি দেয়, তার ফলে একই পদ্ধতি—আরোহী সংজ্ঞা—ব্যবহার করে আমরা সেই
অভিব্যক্তিগুলির অর্থও দিতে পারি।

অভিব্যক্তির অর্থ নির্ধারণ অর্থতত্ত্বের বিষয়। বচনমূলক যুক্তিবিদ্যার
অর্থতত্ত্বে কেন্দ্রীয় ধারণাটি হলো কোনো !!a{valuation}-এ পরিতৃপ্তি।
!!^a{valuation}~$\pAssign{v}$ !!{propositional variable}s-গুলিকে
$\True$, $\False$ সত্যমান আরোপ করে। যেকোনো !!{valuation} যেকোনো
!!{formula}~$!A$-এর জন্য একটি সত্যমান $\pValue{v}(!A)$ নির্ধারণ করে।
কোনো !!^a{formula} যদি !!a{valuation}~$\pAssign{v}$-এ
$\pValue{v}(!A) = \True$ হয়, তবে সেটি সেখানে পরিতৃপ্ত; আমরা এটি
$\pSat{v}{!A}$ লিখি। $!A$-এর গঠনের উপর আরোহের সাহায্যেও এই সম্পর্কটি
সংজ্ঞায়িত করা যায়: যৌক্তিক সংযোজকগুলির সত্য-অপেক্ষক ব্যবহার করে, যেমন
$!A$ ও~$!B$-এর পরিতৃপ্তি (বা অপরিতৃপ্তি) দিয়ে $!A \land !B$-এর
পরিতৃপ্তি সংজ্ঞায়িত করা যায়।

সূত্রগুলির জন্য পরিতৃপ্তি-সম্পর্ক $\pSat{v}{!A}$-এর ভিত্তিতে আমরা
সর্বতঃসত্যতা, অর্থগত অনুসিদ্ধান্ত ও পরিতৃপ্তিযোগ্যতার মৌলিক অর্থগত ধারণা
সংজ্ঞায়িত করতে পারি। কোনো !!^a{formula} সর্বতঃসত্য,
$\Entails !A$, যদি প্রত্যেক !!{valuation} সেটিকে পরিতৃপ্ত করে; অর্থাৎ
যেকোনো~$\pAssign{v}$-এর জন্য $\pValue{v}(!A) = \True$। কোনো
!!{formula}s-এর সেট থেকে সেটি অর্থগতভাবে অনুসৃত হয়,
$\Gamma \Entails !A$, যদি~$\Gamma$-র সব !!{formula}s-কে পরিতৃপ্ত করে
এমন প্রত্যেক !!{valuation}~$!A$-কেও পরিতৃপ্ত করে। আর কোনো
!!{formula}s-এর সেট পরিতৃপ্তিযোগ্য, যদি কোনো !!{valuation} সেটির সব
!!{formula}s-কে একই সঙ্গে পরিতৃপ্ত করে। যেহেতু !!{formula}s আরোহীভাবে
সংজ্ঞায়িত, এবং !!{formula}s-এর গঠনের উপর আরোহের সাহায্যে পরিতৃপ্তিও সংজ্ঞায়িত,
তাই অর্থতত্ত্বের ধর্মগুলি প্রমাণ করতে ও সংজ্ঞায়িত অর্থগত ধারণাগুলির মধ্যে
সম্পর্ক স্থাপন করতে আমরা আরোহ ব্যবহার করতে পারি।


\end{document}
